\section{The Code, the Standards and their enforcement}

\begin{qfdefinition}[The professional conduct programme: who enforces the Code and Standards, how investigations start, the disciplinary review process, the standard of proof and the sanctions]\label{def:professional-conduct-programme}
\qffield{Statement} \emph{Structure.} The CFA Institute Board of Governors oversees the Professional Conduct staff, which investigates, and the Disciplinary Review Committee (DRC), a volunteer committee of members whose panels decide contested cases and review agreed ones. The Bylaws and the Rules of Procedure form the enforcement framework \citep[p.~ix]{standards-practice-handbook-2024}; \citep[Rules 1.1--1.2]{conduct-rules-of-procedure-2022}.

\emph{How investigations start} (Handbook, pp.~ix--x; Rule~2.1): self-disclosure on the annual Professional Conduct Statement of litigation, criminal investigations or written complaints; written complaints from others; public sources such as press and regulatory notices; and, for exams, monitoring of testing centres and social media.

\emph{Process} (Rules 2--5):
\begin{itemize}
\item Staff investigate and may close the case, issue a cautionary letter (Handbook, p.~x), or send a \emph{Statement of Allegations} and then a \emph{Statement of Charges} with findings, violations and a proposed sanction, if a violation is more likely than not (Rules 1.2, 2.5--2.7).
\item The member must answer within 21 days, admitting or denying each finding and accepting or rejecting the sanction (Rule~2.8). A matter can also be settled (Rule~2.9).
\item An accepted charge or settlement with a public sanction goes to a \emph{Review Panel}, which may accept it or impose a lesser sanction but not a greater one (Rules 3.1--3.2). A rejected charge goes to a \emph{Hearing Panel} of three to five DRC members, which decides on the preponderance of the evidence, with the burden on Professional Conduct, and may impose a lesser, equal or greater sanction (Rules 4.2, 4.6--4.7).
\item Suspensions, revocations and prohibitions can be appealed within 28 days; the Appeal Panel looks for clear and material error or an unfair sanction (Rules 5.1, 5.3).
\end{itemize}
\emph{Sanctions} range from a private reprimand to public censure, timed suspension, and revocation of membership and of the right to use the designation; candidates can be suspended from or barred from the exam programme (Rule~1.6). Summary suspension can follow a conviction for a crime punishable by more than a year, a regulatory bar of two years or more, or failure to cooperate (Rule~6.1).

\qffield{Practice} \emph{Path of a case} (original). A charterholder, Dev Arora, discloses on his annual statement that his former employer has filed a complaint about client solicitation. Staff obtain the complaint and client e-mails and conclude that he solicited clients before resigning (Standard~IV(A)). The Statement of Charges proposes a one-year suspension. Arora rejects it, so a Hearing Panel decides on the evidence; it finds the violation more likely than not and imposes censure instead. Arora cannot appeal the censure: Rule~5.1 allows a covered person to appeal only suspensions, revocations and prohibitions.

\qffield{Conditions} Rules of Procedure for conduct related to the profession, as amended and restated 1 January 2022; exam-related matters have separate rules. The Handbook's preface describes the programme as of 2024.

\qffield{Pitfalls} Confusing CFA Institute sanctions with regulatory or criminal penalties: the programme governs membership and the designation, not legal liability. Assuming that silence on a Statement of Charges halts the process: failing to respond counts as acceptance.

\qffield{See also} \hyperref[def:code-of-ethics-standards]{the Code and Standards}; \hyperref[rule:standard-member-responsibilities]{Standard VII}.

\end{qfdefinition}

\begin{qfdefinition}[The Code of Ethics and the Standards of Professional Conduct: six components of the Code, seven Standards and their sub-sections, and the responsibilities each imposes]\label{def:code-of-ethics-standards}
\qffield{Statement} \emph{The Code of Ethics} (six components) requires members and candidates to \citep[p.~xviii]{standards-practice-handbook-2024}:
\begin{enumerate}
\item act with integrity, competence, diligence and respect, and ethically, with the public, clients, prospective clients, employers, employees, colleagues and other market participants;
\item place the integrity of the profession and the interests of clients above their own;
\item use reasonable care and independent professional judgement in analysis, recommendations, investment actions and other professional activities;
\item practise, and encourage others to practise, in a professional and ethical manner that reflects credit on themselves and the profession;
\item promote the integrity and viability of the global capital markets for the ultimate benefit of society;
\item maintain and improve their professional competence and strive to improve that of others.
\end{enumerate}

\emph{The Standards} (pp.~xix--xxii; responsibilities paraphrased):
\begin{itemize}
\item \emph{I Professionalism}: (A) know and comply with applicable law, follow the stricter of law and the Standards, and dissociate from violations; (B) keep independence and objectivity, and refuse gifts that could compromise them; (C) make no misrepresentation; (D) commit no dishonesty, fraud or deceit, or act that reflects adversely on professional reputation, integrity or competence; (E) act with and maintain competence.
\item \emph{II Integrity of capital markets}: (A) do not act, or cause others to act, on material nonpublic information; (B) do not distort prices or inflate volume to mislead.
\item \emph{III Duties to clients}: (A) loyalty, prudence and care; (B) fair dealing; (C) suitability in advisory relationships and consistency with a mandate; (D) fair, accurate and complete performance presentation; (E) confidentiality, except for illegal activities, legal requirements or client permission.
\item \emph{IV Duties to employers}: (A) loyalty; (B) written consent of all parties for additional compensation that could conflict; (C) reasonable efforts by supervisors to secure compliance.
\item \emph{V Investment analysis, recommendations and actions}: (A) diligence and a reasonable, adequate basis; (B) disclosure of services and costs, of the investment process and its changes, of risks and limitations, of the important factors, and separation of fact from opinion; (C) records supporting analysis and communications.
\item \emph{VI Conflicts of interest}: (A) avoid conflicts or disclose them prominently in plain language; (B) client and employer transactions before personal ones; (C) disclosure of referral fees.
\item \emph{VII Responsibilities as a member or candidate}: (A) no conduct that compromises the integrity of CFA Institute programmes; (B) no misrepresentation of membership, the designation or candidacy.
\end{itemize}
The 2024 revision added I(E), the costs-and-services disclosure in V(B) and ``avoid'' in VI(A) (pp.~vii--ix).

\qffield{Practice} \emph{Mapping duties to parties} (original). A research analyst at a bank owes: the market, not to trade on a client's pending takeover (II(A)); the bank's clients, recommendations with a reasonable basis distributed fairly (V(A), III(B)); the bank, loyalty and disclosure of outside paid work (IV(A)--(B)); and readers, disclosure of her own holdings in the stock (VI(A)). One act, such as buying the stock before her upgrade is published, can breach several standards at once (VI(B), III(B)).

\qffield{Conditions} 12th edition of the Handbook (2024); the Standards took effect in revised form on 1 January 2024.

\qffield{Pitfalls} Learning the Standards as isolated rules: the Handbook asks that they be read as one ``interwoven tapestry'' (p.~v). Treating the Code's aspirational components as unenforceable: the Standards implement them.

\qffield{See also} \hyperref[def:professional-conduct-programme]{enforcement}; \hyperref[rule:standard-professionalism]{Standard I} to \hyperref[rule:standard-member-responsibilities]{Standard VII}.

\end{qfdefinition}
